\section{Motivation}
\label{sec:motivation}

\subsection{Challenges}
When a small local model drives the investigation, three challenges arise. We measured each on the baseline configurations of our studies (\S\ref{sec:setup}), in which the model chooses its own probes.

\stepnum{1} \textbf{Probing without converging.} Qwen2.5-7B selecting its own probes verifies only 0.083 to 0.125 of cases. It runs about seven distinct probes per episode, spends 9.4 to 9.5 of its 10 cost units, and issues a verdict in only 6 or 7 of 72 episodes; three quarters of its episodes end with the tool budget exhausted. Replacing its choices with uniformly random ones raises completion by $+0.556$ (\S\ref{sec:rq-rank}), so its own choices are worse than chance.

\stepnum{2} \textbf{Never deciding to stop.} In the decomposition study, under our prompt and serving configuration, Llama-3.1-8B requests another probe in every one of its 3{,}389 decisions, including those in which the ledger's posterior exceeds 0.999 and the prompt states that no legal probe is left. Its outputs are valid JSON throughout; it simply never finishes. Better probing cannot help a model that never states a verdict.

\stepnum{3} \textbf{Verdicts without support.} Larger models conclude, but not always on evidence. Without a ledger, Qwen2.5-32B and 72B call 13 of 51 and 20 of 54 actionable incidents benign among those they give a verdict on, and cite invalid evidence in 22\% and 31\% of their citations (\S\ref{sec:security}). In triage, a verdict that dismisses a real attack is the costliest error, because the incident is closed without a response.

\subsection{Our Solution}
\hesp{} answers each challenge with a component outside the model. \stepnum{1} It ranks every legal probe by expected information gain per cost, from likelihood tables counted in LLM-free development runs rather than elicited from the model, which cannot supply them (elicited tables leave completion at 0.069 to 0.375, \S\ref{sec:rq-help}). \stepnum{2} Before each planner call it applies an explicit acceptance rule and concludes on its own once the rule is met (Appendix~\ref{app:loop} gives the exact order). \stepnum{3} It accepts only verdicts that cite a current observation supporting the claimed cause, and a write-ahead journal lets an audit replay every step.
